\chapter{Discussione}
\label{chap:discussione}

\section{Sintesi dei risultati}
\label{sec:discussione_sintesi}

Il filo che attraversa il caso studio del capitolo precedente si può ricondurre alla tesi formulata nell'Introduzione: una predizione di rischio, per quanto accurata, non è azionabile in mancanza di una spiegazione che indichi su quali fattori intervenire e con quale grado di fiducia. Il confronto condotto sulle stesse istanze mostra che SHAP, LIME e i controfattuali non costituiscono alternative in competizione ma rispondono a tre domande distinte, che insieme compongono un triangolo diagnostico al servizio del responsabile HR. In assenza di tale triangolazione, l'utente si troverebbe nella condizione poco funzionale descritta dalla letteratura sulla sfiducia verso gli algoritmi osservati sbagliare anche una sola volta~\cite{dietvorst2015algorithm}.

La prima dimensione, coperta da SHAP, risponde alla domanda \emph{``perché il modello ha deciso così?''}. Come emerso nella Sezione~\ref{sec:shap}, il \texttt{TreeExplainer} produce un'attribuzione additiva, deterministica e coerente con la decisione effettiva del modello, sia a livello globale (Figura~\ref{fig:shap_summary}) sia a livello della singola istanza (Figura~\ref{fig:shap_wf_214}). La seconda dimensione, coperta da LIME, risponde alla domanda \emph{``quali regole semplici approssimano la decisione?''}: il surrogato lineare locale con feature discretizzate (Sezione~\ref{sec:lime}, Figura~\ref{fig:lime_local_214}) traduce la decisione del modello in soglie e condizioni più vicine al linguaggio di un utente non tecnico, a prezzo di una dipendenza dal campionamento e di una fedeltà puramente locale. La terza dimensione, coperta dai controfattuali Alibi, risponde alla domanda operativamente più incisiva per un HR manager, \emph{``cosa dovrebbe cambiare perché la predizione cambi?''}: l'esito non è un ranking di fattori ma un insieme di modifiche minime alle feature del dipendente capaci di spostare la predizione al di là della soglia di decisione (Sezione~\ref{sec:alibi}, Figure~\ref{fig:cf_214} e~\ref{fig:cf_223}).

La complementarità delle tre prospettive diventa evidente leggendo trasversalmente i tre casi selezionati nelle sezioni precedenti. Per il Dipendente~214, vero positivo ad alta probabilità, i tre framework convergono: SHAP e LIME individuano gli stessi fattori dominanti (straordinari, soddisfazione e retribuzione) e il controfattuale conferma che agendo proprio su quelle leve la predizione scende al di sotto della soglia. In questo caso la spiegazione ha valore di \emph{supporto all'intervento}, perché indica contemporaneamente dove guardare e su cosa agire. Per il Dipendente~223, falso positivo ad alta probabilità, l'allineamento dei tre framework svolge invece una funzione diagnostica sul modello: SHAP e LIME mostrano che il rischio stimato è dominato da caratteristiche riconducibili al profilo giovane e da poco assunto, e il controfattuale rivela che per abbassare la predizione l'ottimizzatore deve intervenire su variabili strutturalmente immutabili come l'età e l'anzianità. La convergenza dei tre metodi segnala che la predizione, pur essendo spiegabile, poggia su un pattern problematico e dunque non andrebbe usata come base di azione. Per il Dipendente~240, falso negativo, SHAP e LIME producono spiegazioni coerenti di permanenza: nessuno dei due metodi, per costruzione, può recuperare informazione assente dalle feature osservate. Qui il limite che le spiegazioni locali rendono visibile è del dataset, non del metodo esplicativo.

Tirando le fila, il caso studio non elegge un ``vincitore'' tra i tre framework. Mostra piuttosto che il valore per il responsabile HR non risiede nella singola probabilità restituita dal modello ma nell'ecosistema che la accompagna: un'attribuzione fondata per la diagnosi, una rappresentazione a regole per la comunicazione, un controfattuale per la prescrizione. In questa chiave, il modello XAI-assistito si configura come strumento di supporto alla decisione e non di sua sostituzione: una distinzione che la sezione successiva riprende per circoscrivere l'ambito di validità delle conclusioni.

\section{Limiti del caso studio}
\label{sec:discussione_limiti}

La lettura dei risultati proposta nella Sezione~\ref{sec:discussione_sintesi} acquista significato solo se accompagnata da una dichiarazione esplicita dei limiti entro cui vale. Tali limiti riguardano la natura del dataset, la sua dimensione, l'insieme delle feature disponibili e le proprietà intrinseche dei tre metodi esplicativi adottati.

Il primo limite riguarda la natura del dataset utilizzato. L'\emph{IBM HR Employee Attrition} è un dataset didattico, costruito a scopo di benchmark e non derivato da un'organizzazione reale: la distribuzione delle feature e la relazione tra queste e il target rispecchiano un'impostazione pensata per la ricerca e la formazione, non le specificità di un contesto aziendale operativo. Di conseguenza, ogni conclusione sulla \emph{performance predittiva} del modello o sulla \emph{composizione dei driver} osservati non può essere trasferita senza cautela a un'organizzazione concreta, pena la confusione tra una proprietà del dataset e una proprietà del fenomeno \cite{IBMHRDataset}.

Il secondo limite è dimensionale. I numeri richiamati nella Sezione~\ref{sec:dataset_setup} --- popolazione ridotta, sbilanciamento marcato, test set che contiene una frazione contenuta di eventi positivi --- hanno due conseguenze concrete. Sulle metriche aggregate di prestazione introducono una varianza campionaria elevata, che rende le stime puntuali di precision, recall e F1 indicative più che stabili. Sulle spiegazioni locali incide invece la rarità dei positivi: le narrazioni costruite sui tre dipendenti selezionati sono plausibili e coerenti, ma la loro \emph{rappresentatività} statistica rispetto alla classe positiva resta limitata proprio dalla dimensione di quest'ultima. L'analisi qualitativa è dunque robusta come illustrazione del comportamento dei tre framework, non come base per inferenze quantitative sulla popolazione.

Il terzo limite riguarda l'insieme delle feature. Il dataset copre dimensioni anagrafiche, retributive, di ruolo, di tenure e di soddisfazione, ma non include variabili di natura organizzativa e relazionale --- clima interno, stile di leadership, offerte esterne ricevute dal dipendente, rete informale di relazioni, eventi personali --- che la letteratura identifica come determinanti spesso decisivi del turnover. Il caso del Dipendente~240, falso negativo con un profilo di permanenza ``pulito'' su tutte le feature osservate, è l'illustrazione più diretta di questo limite: nessun metodo XAI, per quanto sofisticato, può recuperare informazione non presente nei dati di input. Questo vincolo non è specifico del caso studio ma struttura in modo più generale la People Analytics fondata esclusivamente su dati strutturati \cite{Isson_Harriott_2016}.

Il quarto limite è di natura metodologica e riguarda i tre framework stessi, in coerenza con quanto già osservato nel capitolo precedente. La collinearità rilevata tra le variabili di tenure (Sezione~\ref{sec:dataset_setup}) fa sì che SHAP e LIME distribuiscano l'importanza tra proxy dello stesso costrutto, rendendo più difficile isolare una singola leva di intervento. LIME, in quanto basato su campionamento perturbativo, presenta una variabilità residua anche quando il seed è fissato, e produce un surrogato lineare la cui fedeltà è garantita solo localmente. I controfattuali, infine, possono proporre modifiche su feature strutturalmente immutabili (età, anzianità, distanza da casa), come emerso per il Dipendente~223 nella Sezione~\ref{sec:alibi}: in quel caso la spiegazione mantiene un valore diagnostico sul modello, ma perde valore prescrittivo sull'azione HR.

Presi insieme, questi limiti non invalidano il contributo illustrativo del caso studio ma ne circoscrivono l'ambito: le conclusioni tratte valgono come dimostrazione della complementarità dei tre metodi su un benchmark noto, non come stima delle loro prestazioni in contesti organizzativi reali. È su questa consapevolezza che poggia la sezione seguente, dedicata alle implicazioni per la pratica HR.

\section{Implicazioni per la pratica HR}
\label{sec:discussione_implicazioni}

I risultati e i limiti discussi nelle sezioni precedenti acquistano pieno significato solo se tradotti in indicazioni per la pratica del responsabile HR. Questa sezione colloca il triangolo diagnostico all'interno di un workflow operativo, ne discute i rischi legati alla fiducia nel sistema e affronta le implicazioni etiche e normative dell'impiego di modelli predittivi nel contesto del lavoro.

\subsection{Dal segnale all'azione: il workflow XAI-assistito}

La letteratura sulla XAI applicata identifica cinque profili distinti di destinatari delle spiegazioni --- esperti di dominio, soggetti impattati, regolatori, manager e sviluppatori --- ciascuno con esigenze diverse in termini di profondità, formato e scopo della spiegazione~\cite{arrieta2020explainable}. Nel contesto HR, il responsabile del personale appartiene al profilo dell'\emph{esperto di dominio}: non possiede necessariamente competenze di modellazione statistica, ma dispone di conoscenza contestuale sul dipendente che il modello non ha. La spiegazione deve quindi essere \emph{actionable}, cioè capace di colmare il divario tra il segnale computazionale e la decisione operativa~\cite{arrieta2020explainable}.

Il workflow che emerge dal caso studio si articola in tre fasi. Nella prima, il modello segnala i profili a rischio attraverso la soglia business-oriented discussa nella Sezione~\ref{sec:performance}: una soglia calibrata per privilegiare il recall produce un insieme di segnalazioni da sottoporre a triage qualitativo. Nella seconda fase, le spiegazioni attributive (SHAP e LIME) consentono al responsabile HR di comprendere \emph{perché} il modello consideri quel dipendente a rischio, distinguendo tra segnalazioni solide e segnalazioni borderline. Nella terza fase, il controfattuale traduce la diagnosi in prescrizione, indicando su quali leve intervenire per modificare la predizione. Il valore delle spiegazioni locali rispetto alle sole importanze globali è strutturale: le seconde descrivono il modello in aggregato, le prime descrivono il singolo dipendente e abilitano un intervento mirato.

Questa articolazione riflette quanto osservato nelle organizzazioni che utilizzano effettivamente la XAI in produzione: il decisore umano non è un fruitore passivo ma svolge un ruolo di validazione critica, confrontando la spiegazione con la propria conoscenza del contesto e decidendo se confermare, approfondire o annullare la segnalazione del modello~\cite{bhatt2020explainable}.

\subsection{Rischi dell'adozione: algorithm aversion e automation bias}

L'introduzione di un sistema predittivo nel processo decisionale HR espone l'organizzazione a due polarità disfunzionali opposte. Da un lato, la letteratura documenta il fenomeno dell'\emph{algorithm aversion}: la tendenza a rifiutare il supporto algoritmico dopo aver osservato anche un singolo errore, indipendentemente dalla performance complessiva del sistema~\cite{dietvorst2015algorithm}. Dall'altro, studi successivi hanno mostrato il fenomeno speculare dell'\emph{algorithm appreciation}: in assenza di errori visibili, le persone tendono a dare più peso ai consigli algoritmici rispetto a quelli umani, e tale fiducia persiste anche in ambiti soggettivi~\cite{logg2018algorithm}. I due fenomeni coesistono e dipendono dal contesto: l'avversione emerge dopo la visione di un errore, mentre l'apprezzamento è la risposta predefinita in assenza di feedback negativi.

In termini più generali, la tassonomia classica dell'interazione umano-automazione distingue tra \emph{misuse} (sovraffidamento che porta a non rilevare errori del sistema), \emph{disuse} (sottoutilizzazione dovuta a scarsa fiducia) e \emph{abuse} (implementazione del sistema senza considerare le conseguenze sulla prestazione umana)~\cite{parasuraman1997humans}. Il caso studio illustra entrambe le polarità a rischio: un responsabile HR che si affidi acriticamente al punteggio avrebbe attivato un intervento non necessario sul Dipendente~223 (falso positivo); un responsabile che diffidi sistematicamente del modello avrebbe ignorato anche il Dipendente~214 (vero positivo). Il ruolo della XAI è restituire al decisore umano gli elementi per esercitare un giudizio informato, riducendo sia il rischio di misuse sia quello di disuse attraverso la trasparenza della spiegazione.

\subsection{Bias algoritmici e fairness}

Un modello predittivo apprende i pattern presenti nei dati di addestramento, inclusi quelli che riflettono disuguaglianze strutturali. Feature apparentemente neutre possono fungere da \emph{proxy} per caratteristiche protette: la letteratura mostra come variabili quali il quartiere di residenza, la distanza casa-lavoro o la rete di relazioni sociali possano predire con elevata accuratezza l'appartenenza a gruppi demografici, consentendo al modello di discriminare senza mai consultare direttamente i dati sensibili~\cite{barocas2016big}. Inoltre, quando i dati storici riflettono decisioni umane già viziate da pregiudizi, il modello formalizza quelle discriminazioni come regole oggettive, trasformando le disuguaglianze del passato in profezie che si autoavverano~\cite{barocas2016big}.

Nel caso studio, il Dipendente~223 offre una micro-illustrazione di questo rischio: il modello attribuisce un rischio elevato a un profilo giovane con bassa tenure, e il controfattuale rivela che l'unica via per abbassare la predizione è ``invecchiare'' il dipendente. Le spiegazioni locali fungono qui da strumento di \emph{audit}: rendono visibile su quali feature il modello si appoggia e consentono di identificare pattern potenzialmente discriminatori prima che si traducano in azioni. Questa funzione diagnostica è coerente con le strategie di mitigazione documentate nella letteratura sui tool HR, che includono il monitoraggio dell'impatto disparato e la rimozione o riduzione del peso delle feature correlate ad attributi protetti~\cite{raghavan2020mitigating}.

\subsection{Quadro normativo: GDPR e AI Act}

Il quadro normativo europeo impone vincoli specifici all'uso di sistemi decisionali automatizzati nel contesto lavorativo. L'Articolo~22 del GDPR riconosce il diritto dell'interessato a non essere soggetto a decisioni basate unicamente su trattamenti automatizzati che producano effetti giuridici o incidano in modo analogo e significativo sulla persona~\cite{GDPR_Art22}. Nei casi in cui tali decisioni siano consentite, il regolamento prescrive misure di tutela che includono il diritto all'intervento umano, a esprimere la propria opinione e a contestare la decisione~\cite{goodman2017eu}.

Il dibattito in letteratura si concentra sulla portata dell'obbligo di fornire ``informazioni significative sulla logica coinvolta''. Un'interpretazione restrittiva limita tale obbligo a una descrizione ex ante della funzionalità del sistema; un'interpretazione funzionale sostiene invece che, poiché i modelli di machine learning sono deterministici, una spiegazione che non permetta all'interessato di comprendere come i propri dati abbiano influenzato la decisione non può considerarsi significativa~\cite{selbst2017meaningful}. In questa seconda ottica, le spiegazioni controfattuali sono state proposte come la forma più compatibile con il requisito del GDPR, perché indicano le condizioni minime sotto le quali la decisione sarebbe stata diversa senza richiedere la disclosure dell'intero modello~\cite{wachter2017counterfactual}.

Il Regolamento (UE) 2024/1689 (AI Act) introduce un ulteriore livello di disciplina. L'Articolo~6, paragrafo~2, in combinazione con l'Allegato~III, punto~4, classifica esplicitamente i sistemi di intelligenza artificiale utilizzati per il recruitment, la gestione del rapporto di lavoro, le decisioni su promozioni e cessazioni come sistemi \emph{ad alto rischio}~\cite{EU_AI_Act_2024}. Per tali sistemi, l'Articolo~13 impone obblighi di trasparenza --- istruzioni d'uso che includano caratteristiche, limiti di prestazione e misure di supervisione --- e l'Articolo~14 prescrive che siano progettati per essere effettivamente sorvegliati da persone fisiche dotate della competenza, formazione e autorità necessarie per comprenderne i limiti e, se del caso, ignorarne o annullarne l'output~\cite{EU_AI_Act_2024}.

\subsection{Human-in-the-loop come requisito di governance}

Le considerazioni precedenti convergono su un principio trasversale: il modello XAI-assistito non è un sistema decisionale autonomo ma uno strumento che amplifica la capacità analitica del responsabile HR, a condizione che resti inserito in un ciclo di supervisione umana effettiva. Il valore della spiegazione non è un fine in sé, ma il mezzo attraverso cui tale supervisione diventa informata e capace di individuare tanto i segnali affidabili quanto i pattern problematici. In questa prospettiva, la XAI non risolve il problema della fiducia nel modello ma lo trasforma: da una scelta binaria (fidarsi o meno della predizione) a una valutazione graduata, fondata sulla qualità e sulla coerenza della spiegazione che la accompagna.
